\subsection{实直线与实数}

定义 2. 令 R 表示作为有理直线加法群 Q 的完备化的那个拓扑群。R 的元素称为实数；作为拓扑空间，R 称为实直线；作为拓扑群，R 称为实直线的加法群。

我们恒把 Q 与 R 中那个典范同构于它的稠密子群视为同一。在这一约定下，每个有理数都是实数。每个非有理数的实数称为无理数；我们已在第二章 § 3 第 3 目中见过这样的数存在（本章 § 3 第 3 目将用另一种方式证明这一点；另参看 § 2 的习题 2）；因此（第三章 § 2 第 1 目）无理数集 CQ 在 R 中稠密。

我们将证明 Q 的序结构可以延拓到 R 上，使延拓后的序仍与 R 的加法群结构相容：

命题 1. 关系 \(y - x \in \overline{\mathbf{Q}}_+\) 是 \(\mathbf{R}\) 上的一个序，它使 \(\mathbf{R}\) 成为全序集；它与 \(\mathbf{R}\) 的加法群结构相容，并且把序 \(x \leqslant y\) 诱导在 \(\mathbf{Q}\) 上。

我们先证明关系 \(y - x \in \overline{\mathbf{Q}}_+\) 与 \(z - y \in \overline{\mathbf{Q}}_+\) 蕴含 \(z - x \in \overline{\mathbf{Q}}_+\) 。事实上，函数 \(x + y\) 在 \(\mathbf{R} \times \mathbf{R}\) 上连续，因此由 (8) 我们有 \(\overline{\mathbf{Q}}_+ + \overline{\mathbf{Q}}_+ \subset \overline{\mathbf{Q}}_+\) （第一章 § 2 第 1 目，定理 1）。其次，我们要证明关系 \(y - x \in \overline{\mathbf{Q}}_+\) 与 \(x - y \in \overline{\mathbf{Q}}_+\) 蕴含 \(x = y\) ；这就确立了 \(y - x \in \overline{\mathbf{Q}}_+\) 是 \(\mathbf{R}\) 上的一个序。只需证明 \(\overline{\mathbf{Q}}_+ \cap (-\overline{\mathbf{Q}}_+) = \{0\}\) 。而函数 \(x \to x^+\) 与 \(x \to x^-\) 在 \(\mathbf{Q}\) 上一致连续，因此可以连续延拓到 \(\mathbf{R}\) （第二章 § 3 第 6 目，定理 2）；令 \(f\) 与 \(g\) 分别为它们的延拓。由延拓，\(x = f(x) - g(x)\) 对所有 \(x \in \mathbf{R}\) 成立；若 \(x \in \overline{\mathbf{Q}}_+\) 则 \(g(x) = 0\) ，又由于 \(-\overline{\mathbf{Q}}_+\) 是 \(-\mathbf{Q}_+\) 的闭包（这由 \(-x\) 的连续性得出），于是若 \(x \in -\overline{\mathbf{Q}}_+\) 则 \(f(x) = 0\) 。因此若 \(x \in \overline{\mathbf{Q}}_+ \cap (-\overline{\mathbf{Q}}_+)\) ，就有 \(f(x) = g(x) = 0\) ，从而 \(x = 0\) 。

由 (10)，我们有 \(\overline{\mathbf{Q}}_{+} \cup (-\overline{\mathbf{Q}}_{+}) = \mathbf{R}\) ，因而 \(\mathbf{R}\) 被序 \(y - x \in \overline{\mathbf{Q}}_{+}\) 全序化。

此外，由于关系 \(y - x \in \overline{\mathbf{Q}}_+\) 与 \((y + z) - (x + z) \in \overline{\mathbf{Q}}_+\) 等价，序 \(y - x \in \overline{\mathbf{Q}}_+\) 与 \(\mathbb{R}\) 的加法群结构相容。

最后，若 \(x\) 与 \(y\) 属于 \(\mathbf{Q}\) ，则关系 \(y - x \in \overline{\mathbf{Q}}_+\) 与 \(y - x \in \mathbf{Q}_+\) 等价，因此关系 \(y - x \in \overline{\mathbf{Q}}_+\) 把关系 \(x \leqslant y\) 诱导在 \(\mathbf{Q}\) 上。证毕。

关系 \(y - x \in \overline{\mathbf{Q}}_+\) 仍记作 \(x \leqslant y\) 。集合 \(\overline{\mathbf{Q}}_+\) 是所有 \(x \geqslant 0\) （ \(\mathbf{R}\) 中的）组成的集，记作 \(\mathbf{R}_+\) ；它是一个闭集。所有 \(x > 0\) 组成的集记作 \(\mathbf{R}_+^*\) ；它是 \(-\mathbf{R}_+\) 的余集，因而在 \(\mathbf{R}\) 中是开集。

